\chapter{Lizenzen}  % Kapitel % Steht dann über dem Text
\label{chapter:Lizenzen}  % Steht als Text im Inhaltsverzeichnis
\index{Lizenzen} % für das Stichwortverzeichnis

\begin{figure}[h]
\begin{center}
   \includegraphics[width=130pt]{../bilder/gnu-head-sm.jpg}
   \caption{GNU Projekt Logo}
   \label{GPL Lizenz}
\end{center}
\end{figure}
Das Projekt CodeRed unterliegt in seinen Grundzügen der \href{http://de.wikipedia.org/wiki/GPL}{GPL Lizenz} und wurde mit Hilfe von Freien Software Produkten und Entwicklungsumgebungen geschaffen.\\
\begin{figure}[h]
\begin{center}
   \includegraphics[width=90pt]{../bilder/mysql.jpg}
   \includegraphics[width=90pt]{../bilder/rails.png}
   \includegraphics[width=90pt]{../bilder/debian.png}
   \includegraphics[width=90pt]{../bilder/apache.png}
   \caption{GNU Projekt Logo}
   \label{GPL Lizenz}
\end{center}
\end{figure}
Marco Benecke und Jan Neuser sind die Urheber von CodeRed und haben daher alle geltenden Rechte an dem Trouble Ticket System CodeRed.